\section{What Does a Face-Reading System Measure?}
\label{sec:11.8}

\textbf{Unresolved issue.}
Facial movement alone does not reliably identify a person's emotional state, and recognition accuracy varies across cultural and linguistic groups \autocite{barrett2019emotional,elfenbein2002universality}. It remains unresolved whether a face-reading system measures emotion at all or primarily measures distance from the reference population and registration model chosen by its developers.

\textbf{Test.}
Audit the system's registration stage as well as its final classifications. Publish the reference model, or sufficient summary statistics to reconstruct it, and measure each evaluated face's distance from that reference. Report error as a function of this continuous distance and separately compare performance on populations represented and not represented in the training data. The evaluation corpus must record the relevant cultural context and the procedure used to establish each target label.

\textbf{Evidence against the proposal.}
If error does not increase with distance from the reference and remains flat across represented and unrepresented populations, the manuscript's account of the system as measuring deviation from a chosen norm is wrong for that system. If error does increase with distance, the result counts against using the system as evidence of another party's emotional state.

\textbf{Required access or authority.}
The audit requires the registration model, training-corpus composition, evaluation outputs, and documentation of the labeling procedure. Because vendors have little incentive to expose their reference models, an independent auditor must have authority to compel access and publish disaggregated results.
